\section{Energy budgets and the common capacity dual}
\label{sec:energy-framework}

Both compact cases use the same specialization of
Theorem~\ref{thm:homogeneous-certificate}. A complete domain $b$ has
its own original orthogonal energy slots and all applicable rows
\begin{equation}\label{eq:branch-energy}
                  A_be_b\ge c_b,\qquad e_b\ge0.
\end{equation}
Here $A_b\ge0$ is rational, each row has a positive entry, and
$c_b\ge0$ is rational. In the actual matrices $A_b$ records supports
by $0$ and $1$. More generally, let the fixed geometric budget be
\begin{equation}\label{eq:general-energy-budget}
              h_b^Te_b<B_b,\qquad h_b>0,\quad B_b>0.
\end{equation}
The weights $h_b$ and budgets $B_b$ are rational data.
The compact geometry gives $h_b=\one$ and $B_b=1$ by
\eqref{eq:component-energies}. Leaving them as data shows which part
of the formulation is independent of that numerical specialization.

Different complete domains retain separate coordinates. For example,
a two-slot configuration and a three-slot configuration do not share
energy variables merely because both occur in the cover. Within one
domain, however, every use of a given original component refers to
its one slot. This distinction preserves both local applicability and
energy ownership.

\subsection{One normalization and duality theorem}

For each domain introduce nonnegative mass $\tau_b$, scaled energies
$E_b$ and deficit $\delta_b$. Impose
\begin{equation}\label{eq:budget-deficit}
 \begin{gathered}
 A_bE_b\ge c_b\tau_b,\qquad
 \delta_b=B_b\tau_b-h_b^TE_b\ge0,\\
                         \sum_b\delta_b=1.
 \end{gathered}
\end{equation}
A proved local exclusion adds
\begin{equation}\label{eq:excluded-mass}
                              \tau_b=0.
\end{equation}
It then forces $E_b=0$ and $\delta_b=0$, since $h_b>0$.
A fixed positive tolerance is unnecessary: the positive physical
budget deficit itself supplies the normalization.

\begin{theorem}[Budget lift and capacity dual]
\label{thm:budget-lift}\label{prop:general-energy-budget}
\label{thm:energy-normalization}\label{thm:energy-certificate}
\label{prop:capacity-complete}
For a finite complete domain family, the following assertions hold.
\begin{enumerate}[label=\textup{(\roman*)}]
\item System~\eqref{eq:budget-deficit}--\eqref{eq:excluded-mass} is
feasible exactly when a nonexcluded domain admits energies satisfying
\eqref{eq:branch-energy} and~\eqref{eq:general-energy-budget}.
Every actual polytope in the covered dimensions therefore has a lift.
\item Put
\begin{equation}\label{eq:branch-minimum}
 m_b=\min\{h_b^Te:e\ge0,\ A_be\ge c_b\}.
\end{equation}
Then
\begin{equation}\label{eq:branch-duality}
 m_b=\max\{c_b^T\alpha:\alpha\ge0,\ A_b^T\alpha\le h_b\},
\end{equation}
and both optima are attained rationally. The normalized system is
infeasible exactly when every
nonexcluded domain has $m_b\ge B_b$.
\item Consequently nonnegative rational multipliers satisfying
\begin{equation}\label{eq:general-capacity-certificate}
               A_b^T\alpha_b\le h_b,\qquad
               c_b^T\alpha_b\ge B_b
\end{equation}
for each nonexcluded domain give a complete infeasibility certificate
for this linear energy model.
\end{enumerate}
\end{theorem}
\begin{proof}
For an energy vector with $s=B_b-h_b^Te_b>0$, take
$\tau_b=1/s$, $E_b=e_b/s$, $\delta_b=1$ and set all other domains
zero. Conversely, some deficit in a normalized feasible point is
positive. Its mass is positive, so division by that mass gives the
local rows and
\[
           h_b^T(E_b/\tau_b)=B_b-\delta_b/\tau_b<B_b.
\]
Completeness of the geometric cover proves the actual-object claim.
The masses can be arbitrarily large; no artificial upper bound is
imposed on them.

For the optimum, a sufficiently large multiple of $\one$ satisfies
all local rows. Positivity of $h_b$ makes every finite objective
sublevel compact, so the minimum is attained. To see the dual signs
and equality directly, omit the subscript and consider the polyhedral
cone
\[
       K=\{(Ae-u,h^Te+t):e,u,t\ge0\}.
\]
It contains $(c,m)$ but no $(c,m-\varepsilon)$ for
$\varepsilon>0$. A nonzero supporting functional $(-\alpha,\beta)$,
nonnegative on $K$ and zero at $(c,m)$, gives
\[
 \alpha\ge0,\quad\beta\ge0,\quad A^T\alpha\le\beta h,
                  \quad\alpha^Tc=\beta m.
\]
If $\beta=0$, the positive entry in each nonnegative row of $A$
forces $\alpha=0$, a contradiction. Division by $\beta$ gives a
dual optimizer. The reverse inequality follows by multiplying
$Ae\ge c$ by any feasible dual vector. A vertex of the compact
primal optimal face is a rational vertex of the original feasible
polyhedron. The resulting dual optimal face is also a nonempty
rational polyhedron, so has a rational point by
Lemma~\ref{lem:cone-alternative}.

Finally, multiplication of the scaled rows by a capacity certificate
shows
\[
 h_b^TE_b\ge\alpha_b^TA_bE_b
            \ge(c_b^T\alpha_b)\tau_b\ge B_b\tau_b.
\]
Thus every deficit is zero. In the common formulation, take $C_b$ to
be the local cone in~\eqref{eq:budget-deficit}, with zero mass on
excluded domains, $T_b=0$ and $d_b=\delta_b$.
The displayed inequality is precisely the certificate
$-d_b\ge0$ of Theorem~\ref{thm:homogeneous-certificate}.
Part (ii) makes this certificate criterion necessary as well as
sufficient for the stated finite linear model.
\end{proof}

\subsection{Capacities and optimization in the compact applications}

For $h_b=\one$, $B_b=1$, the implemented dual seeks
\begin{equation}\label{eq:capacity-dual}
 \alpha_b\ge0,\qquad A_b^T\alpha_b\le\one,\qquad
              c_b^T\alpha_b\ge\eta,\qquad\eta\ge1.
\end{equation}
There is one capacity inequality for each original energy slot. For
support rows its left side is
$\sum_{\ell:j\in S_\ell}\alpha_{b\ell}$: the total multiplicity with
which that slot is charged. A new local theorem does not give another
copy of the same energy.

\begin{example*}
Suppose $e_I+e_A\ge a$ and $e_I+e_B\ge b$. Weights one on both rows
would count $e_I$ twice. Weights $1/2,1/2$ give instead
\[
           e_I+\tfrac12e_A+\tfrac12e_B\ge\tfrac12(a+b),
\]
a valid lower bound for the total energy because the omitted
half-energies are nonnegative. Capacities test validity of the
combination; its payment tests whether it reaches the budget.
\end{example*}

For one fixed dimension-labelled catalogue define
\begin{equation}\label{eq:threshold-energy-value}
 \Gamma_N=\min_{\substack{b\ \text{not locally excluded}\\ n_b\ge N}}m_b,
 \qquad\min\varnothing=+\infty.
\end{equation}
The normalized LP is infeasible exactly when $\Gamma_N\ge1$.
Raising $N$ removes domains, so this minimum cannot decrease; this
is the energy form of Proposition~\ref{prop:threshold}.
Completeness concerns the supplied linear inequalities, not geometric
realizability of an arbitrary feasible energy vector.

The recorded computation enforces~\eqref{eq:capacity-dual} and
maximizes the sum of local payments, favoring simple saturated
capacity solutions. Exact verification of the resulting rational
point proves the needed comparison. No numerical optimality claim
is used. Compactness and compact arithmeticity change the domains
and rows; their linear assembly is identical.
